\documentclass[11pt,a4paper]{article}

\usepackage[margin=2.5cm]{geometry}
\usepackage[french]{babel}
\usepackage[T1]{fontenc}
\usepackage[utf8]{inputenc}
\usepackage{setspace}
\usepackage{hyperref}

\setstretch{1.2}

\begin{document}

%========================
% Coordonnées du demandeur
%========================
\begin{flushright}
Philippe Thomas Savard\\
4-5354, rue du Menuet\\
Lévis (secteur Charny), Québec\\
Canada\

\[4pt]
Courriel : \texttt{philippethomassavard@exemple.com}\\
Téléphone : \texttt{(418) 487-8798}\

\[8pt]
Lévis, le \vingt cinq mars deux milles vintg-six
\end{flushright}

\vspace{1cm}

\noindent
\textbf{À l'attention de :}\\
Ville de Lévis\\
Service des travaux publics et/ou de la circulation

\vspace{0.5cm}

\begin{center}
\Large\textbf{Requête concernant la sécurité hivernale et le stationnement}\

\[4pt]
\large\textbf{Rue du Menuet et rue du Mélomane – secteur Charny}
\end{center}

\vspace{0.8cm}

%========================
% Introduction
%========================
\section*{Introduction}

Je, soussigné, \textbf{Philippe Thomas Savard}, citoyen de la Ville de Lévis, résident au \textbf{4-5354, rue du Menuet}, secteur Charny, souhaite porter à votre attention une problématique de sécurité liée au stationnement et aux opérations de déneigement et de déglaçage sur une portion de la rue du Menuet et de la rue du Mélomane.

Le secteur concerné se situe \textbf{devant l'immeuble du 5354, rue du Menuet} et s'étend \textbf{jusqu'à l'intersection de la rue du Mélomane et de la rue de la Rapsodie}. À mon avis, cette section du réseau routier présente des dangers importants pour les citoyens, les piétons, les automobilistes et les autres usagers, particulièrement durant la saison froide.

Le présent document a pour but d'exposer la situation, d'en expliquer les conséquences et de formuler une \textbf{demande officielle d'interdiction de stationnement en période hivernale} sur ce tronçon.

%========================
% Exposé de la problématique
%========================
\section*{Exposé de la problématique}

Durant la période hivernale, lorsque les services de déneigement et de déglaçage de la Ville de Lévis sont en opération, la portion de la rue du Menuet située devant le 5354, ainsi que le tronçon menant vers l'intersection de la rue du Mélomane et de la rue de la Rapsodie, est fréquemment occupée par des véhicules stationnés des deux côtés de la chaussée.

Cette situation entraîne plusieurs conséquences :

\begin{itemize}
  \item \textbf{Réduction importante de la largeur de la chaussée} : la présence de véhicules stationnés des deux côtés de la rue rétrécit considérablement l'espace disponible pour la circulation et pour le passage de la machinerie de déneigement et d'épandage d'abrasif.
  \item \textbf{Impossibilité d'épandre l'abrasif de manière adéquate} : afin de protéger les véhicules stationnés, les opérateurs des équipements d'épandage doivent interrompre ou limiter l'application d'abrasif sur cette portion de rue, ce qui laisse la chaussée sans traitement suffisant.
  \item \textbf{Formation de glace et chaussée très glissante} : en raison de l'absence ou de l'insuffisance d'abrasif, la chaussée demeure souvent glacée et dangereuse sur une grande partie de la saison froide.
  \item \textbf{Danger pour les piétons} : les allées piétonnières donnant accès à la rue devant le 5354, rue du Menuet, deviennent difficiles et risquées à emprunter. Les piétons doivent parfois marcher sur une portion de la chaussée, déjà glissante, pour se déplacer en sécurité relative.
  \item \textbf{Risque accru pour les automobilistes et les usagers} : la combinaison d'une chaussée rétrécie, de la glace persistante et de la circulation locale crée un environnement routier hasardeux, particulièrement lors des périodes de froid intense et de redoux.
\end{itemize}

De manière générale, la \textbf{chaussée et les accès piétonniers} devant le 5354, rue du Menuet, et jusqu'à l'intersection de la rue du Mélomane et de la rue de la Rapsodie, se trouvent dans un état jugé dangereux une grande partie de la saison hivernale, principalement en raison du stationnement sur rue qui empêche un déneigement et un déglaçage efficaces.

%========================
% Demande formulée à la Ville
%========================
\section*{Demande formulée à la Ville de Lévis}

Par le présent document, je formule respectueusement la demande suivante :

\begin{quote}
\textbf{Que la Ville de Lévis analyse la possibilité d'interdire le stationnement en période hivernale, en tout temps, sur la portion de la rue du Menuet située devant le 5354, ainsi que sur la section de la rue du Mélomane jusqu'à l'intersection de la rue du Mélomane et de la rue de la Rapsodie durant la saison froide}
\end{quote}

L'objectif de cette demande est de :

\begin{itemize}
  \item permettre aux équipes de déneigement et de déglaçage d'effectuer leur travail de manière complète et sécuritaire ;
  \item assurer une application adéquate d'abrasif sur la chaussée ;
  \item réduire les risques de chute et d'accident pour les piétons ;
  \item améliorer la sécurité routière pour l'ensemble des usagers de cette portion de rue.
\end{itemize}

Je suis conscient que des règles de stationnement spécifiques existent déjà lors des opérations de déneigement. Toutefois, l'expérience des dernières saisons hivernales montre que la situation demeure problématique et dangereuse. C'est pourquoi je sollicite une \textbf{interdiction de stationnement hivernal en tout temps} sur ce tronçon précis.

%========================
% Documentation complémentaire
%========================
\section*{Documentation complémentaire}

Afin de documenter plus en détail la situation, j'ai créé un dépôt public sur la plateforme GitHub, où sont rassemblés :

\begin{itemize}
  \item des photographies illustrant l'état de la chaussée et du stationnement en période hivernale ;
  \item le texte de la présente requête ;
  \item des informations complémentaires sur le secteur concerné.
\end{itemize}

Le dépôt est accessible à l'adresse suivante :

\begin{center}
\url{https://github.com/2racinede4carreunivers-dev/requete_d-galcage_menuet_melomane_2026_citoyen.git}
\end{center}

%========================
% Conclusion
%========================
\section*{Conclusion}

En résumé, je demande à la Ville de Lévis d'examiner attentivement la situation décrite et de considérer l'instauration d'une \textbf{interdiction de stationnement en saison froide} sur la portion de la rue du Menuet devant le 5354, ainsi que sur la section de la rue du Mélomane jusqu'à l'intersection avec la rue de la Rapsodie.

Je demeure disponible pour toute précision supplémentaire, pour fournir d'autres documents ou pour participer à toute démarche de consultation jugée pertinente.

\vspace{0.8cm}

Je vous remercie de l'attention portée à la présente requête. 

\vspace{1.2cm}

\noindent
Veuillez agréer l'expression de mes salutations distinguées.

\vspace{1.5cm}

\noindent
\textbf{Philippe Thomas Savard}\\
Citoyen de la Ville de Lévis, secteur Charny

\end{document}
